% ============================================================================
% sec3_imaging.tex — Filmed glass: imaging, capture and panel (§3, W1 rev2)
% Target: ~545 words (Lead: must fit the 5.12 p body budget; trim in-section).
% Numbers only from plan 附B (B01-B14, B06-B07); Table I lives in supplementary S2.
% ============================================================================
\section{Filmed Glass: Imaging, Capture and Panel}
\label{sec:imaging}

\subsection{What the film does to the two layers}
\label{sec:film}

For a sensor measurement through a filmed pane,
$I(x)=\tau\,T(x)+\rho\,R(x)$, where $\tau$ is the film-attenuated transmission, $\rho$
the reflection coefficient, $T$ the scene and $R$ the reflected environment. Standard
SIRR implicitly assumes that $\rho R$ is bright and $\tau T$ well exposed; window film
challenges both. Since the film is laminated on the scene side and the camera on the
bare side, transmitted light passes \emph{through} the film bulk and is scattered,
absorbed and dyed, whereas the reflection of the behind-camera environment comes off the
bare \emph{surface} and enters no film. The predicted asymmetry in sharpness and chroma
is what we observe. With $\mathrm{sharp}(X)=\mathbb{E}|\mathrm{HP}(X)|/\mathbb{E}|X|$
(HP: unit-variance high-pass) averaged per dataset, filmed glass has
$\mathrm{sharp}(T)=0.027$ against $0.041$--$0.062$ on SIR$^2$-Wild, SIR$^2$-Obj and
Real20, i.e.\ its transmission is about $2.3\times$ softer; the
reflection-to-transmission sharpness ratio is $6.76$ against $2.51$--$3.32$ ($2.0\times$
sharper relative to the layer beneath); and the transmission chroma is $0.219$ against
$0.322$--$0.393$. Mean brightness, mean reflection magnitude and the high-frequency
\emph{fraction} of the reflection are, by contrast, indistinguishable from Real20 when
measured identically on the four datasets: the asymmetry, not brightness, separates the
regimes (full table in the supplementary material).

The film also changes \emph{which} pixels are damaged. Sorting scene pixels by
ground-truth luminance, the reflection magnitude $|I-T|$ is $23.1$ gray levels in the
darkest quartile against $14.1$ in the brightest (the four quartiles give $23.1$, $25.7$,
$17.7$, $14.1$): the two layers are degraded \emph{in the same places}, the opposite of
the usual SIRR setting, where a bright reflection sits on a well-exposed scene.

\subsection{Capture protocol and the ground-truth hazard}
\label{sec:protocol}

All scenes are shot on a phone through an automotive window film, handheld, with
exposure and white balance locked. For each scene we record a short clip with reflection
and then cover the glass with dark cloth to obtain $T$: one $T$ per scene, many $I$;
inputs are extracted at 30\,fps, cropped to the glass area, and a pair is kept only if
the phase-correlation shift is $\le2$\,px and the per-channel gain deviation
$\le0.15$. Covering the glass is not neutral, and on a pool of $82$ confirmed pairs we
report three separate measurements rather than a single causal chain: the transmission
reference is $3.7\%$ darker inside the hand-annotated occluder region than outside
($T_{\mathrm{in}}/T_{\mathrm{out}}=0.963$); the reflected content inside the region is
lower than outside by $16.8\%$ ($r_{\mathrm{in}}/r_{\mathrm{out}}=0.832$); and the
input-to-target ratio is $I/T|_{\mathrm{in}}\approx0.887$ against $1.045$ outside. The
reflection contributes a median $4.5\%$ over the target outside the region ($33\%$ in the
strongest decile). Exposure must be locked: with auto-exposure active we measured
$\mathrm{median}(I/T)=0.99$ and a residual layer correlated with image brightness at
$r=0.32$; after locking, $r=-0.01$. Because $T$ is a separate frame, alignment is a
second hazard: on one batch $88.7\%$ of the pairs were shifted by more than $3$\,px
(median $8.2$\,px). Per-pair translation alignment is a single-seed screening result on
this pool (aligned pool $21.43$\,dB; discarding the $99$ worst pairs gives $+0.80$\,dB)
and is not used for any main-table claim.

\subsection{Development panel and metrics}
\label{sec:panel}

The panel contains 15 scenes ($1280{\times}832$) shot through the filmed pane. For every
scene we hand-annotated the occluder region $\Om$ once (median $4.6\%$ of the image,
range $2.9$--$9.2\%$, under a fixed intensity rule); it marks where the reflection path
was blocked, and it is an annotation, not a physical ground truth. We report PSNR, SSIM,
HF-PSNR, gradient cosine, the Laplacian-variance ratio
$\mathrm{LapVar}(\That)/\mathrm{LapVar}(T)$ and region-restricted PSNR inside/outside
$\Om$, plus the residual-reflection coefficient
$\beta=\langle T-\That,\,I-T\rangle/\lVert I-T\rVert^{2}$ ($-1$: nothing removed, $0$:
perfect) and the per-image gain
$\alpha=\mathrm{median}_i(\overline{\That_i}/\overline{T_i})$. With $n=15$ scenes the
minimum detectable effect of the paired panel is $\approx0.26$\,dB, and the floor for
brightness-only effects is $1.19$--$1.31$\,dB. The panel was used while developing the
annotation protocol and the pilot models, so we call it a \emph{development panel} and
treat its numbers as development-set evidence rather than confirmation.
